\kbvsechead{ನಾಮ ಪ್ರಕೃತಿ :}

\indent ಯಾವುದೇ ವ್ಯಕ್ತಿಯ ಹೆಸರನ್ನಾಗಲೀ, ವಸ್ತುವಿನ ಹೆಸರನ್ನಾಗಲೀ, ಸ್ಥಳದ ಹೆಸರನ್ನಾಗಲೀ, ಸಸ್ಯದ ಹೆಸರನ್ನಾಗಲೀ, ಪಶುಪಕ್ಷಿಯ ಹೆಸರನ್ನಾಗಲೀ ಸೂಚಿಸುವ ಪದಕ್ಕೆ “ನಾಮ ಪ್ರಕೃತಿ” ಎಂದು ಹೆಸರು.\par

\noindent\hspace*{1.2em}\kbvlabel{ಉದಾ}\hspace{0.6em}ಲಕ್ಷ್ಮಣ, ಮಾಲಿನಿ, ಕಬ್ಬಿಣ, ಕಲ್ಲು, ಕಾಡು, ಊರು, ಗುಲಾಬಿ, ಬಾಳೆ, ಎಮ್ಮೆ, ಪಾರಿವಾಳ ಇವು ನಾಮ ಪ್ರಕೃತಿಗಳು.\par

\kbvsechead{ಪ್ರತ್ಯಯ :}

\noindent\hspace*{2.2em}ಹೊಲವು ಊರಿಗೆ ಸಮೀಪದಲ್ಲಿದೆ.\par

\noindent\hspace*{2.2em}ಹೊಲದಲ್ಲಿ ಕಸಕಡ್ಡಿ ತುಂಬಿದೆ.\par

\noindent\hspace*{2.2em}ರೈತನು ಹೊಲವನ್ನು ಉಳುತ್ತಿದ್ದಾನೆ.\par

\noindent\hspace*{2.2em}ಹೊಲದ ಅಂಚಿನಲ್ಲಿ ಹುತ್ತವಿದೆ.\par

\indent ಈ ವಾಕ್ಯಗಳಲ್ಲಿ ದಪ್ಪ ಅಕ್ಷರಗಳಲ್ಲಿರುವ ಹೊಲವು, ಹೊಲದಲ್ಲಿ, ಹೊಲ\par

\indent ವನ್ನು, ಹೊಲದ-- ಪದಗಳನ್ನು ಗಮನಿಸಿ. ಇಲ್ಲಿ “ಹೊಲ” ಎಂಬ ನಾಮ ಪ್ರಕೃತಿಯನ್ನು ಸಂದರ್ಭೋಚಿತ ರೀತಿಯಲ್ಲಿ ಬಳಸಿರುವುದು ಕಂಡುಬರುತ್ತದೆ.\par

\noindent\hspace*{2.2em}ಈಗ ಆ ಪದಗಳನ್ನು ಬಿಡಿಸಿ ನೋಡೋಣ.\par

\begin{kbvderiv}
\vspace{4pt plus 3pt minus 1pt}
\begin{center}\begin{adjustbox}{max width=\textwidth}
{\renewcommand{\arraystretch}{1.22}\setlength{\tabcolsep}{4pt}
\begin{tabular}{l c l}
ಹೊಲ + ಉ & {\color{kbvgrey}$=$} & ಹೊಲವು (ವ್ + ಉ = ವು) \\
ಹೊಲ + ಅಲ್ಲಿ & {\color{kbvgrey}$=$} & \kbvresult{ಹೊಲದಲ್ಲಿ} \\
ಹೊಲ + ಅನ್ನು & {\color{kbvgrey}$=$} & \kbvresult{ಹೊಲವನ್ನು} \\
ಹೊಲ + ದ & {\color{kbvgrey}$=$} & \kbvresult{ಹೊಲದ} \\
\end{tabular}}
\end{adjustbox}\end{center}
\vspace{4pt plus 3pt minus 1pt}
\end{kbvderiv}

\indent ಮೇಲಿನ ನಾಲ್ಕು ಸಂದರ್ಭಗಳಲ್ಲಿಯೂ ಹೊಲ ಎಂಬ ನಾಮ ಪ್ರಕೃತಿ ಯೊಂದಿಗೆ ಕ್ರಮವಾಗಿ ಉ, ಅಲ್ಲಿ, ಅನ್ನು, ಆ ಸೇರಿರುವುದನ್ನು ಕಾಣಬಹುದು ; ಇವೇ “ಪ್ರತ್ಯಯಗಳು” ಎಂದು.\par

\begin{kbvrule}
\noindent\kbvbullet\hspace{0.5em}ಪ್ರಕೃತಿಯೊಂದಿಗೆ ಪ್ರತ್ಯಯ ಸೇರಿ ಪದವಾಗುವುದು. ಆದುದ ರಿಂದ ಹೊಲವು, ಹೊಲದಲ್ಲಿ, ಹೊಲವನ್ನು, ಹೊಲದ ಎಂಬ ಪದಗಳು ನಾಮ ಪದಗಳಾಗಿರುವುವು.
\end{kbvrule}

\begin{kbvrule}
\noindent{\kbvsecfont ವ್ಯಾಖ್ಯೆ :}\hspace{0.5em}ನಾಮ ಪ್ರಕೃತಿಗಳ ಮುಂದೆ ಪ್ರತ್ಯಯಗಳು ಸೇರಿ ನಾಮ ಪದಗಳೆನಿಸುವುವು.
\end{kbvrule}
